\section{The heptagon residual has a five-point transversal}\label{sec:deletion}
Write $B=M_7(5,6)$. Call a point a \emph{right root} if it starts no
three-cap in the specified ambient set. Every triple starting at a
right root is a cup, so it can be prepended to any cup to its right.

\begin{lemma}[Endpoint shift]\label{lem:root}
Every endpoint of a six-cup in a set counted by $D_7$ is a right root.
\end{lemma}
\begin{proof}
Let a six-cup end with $u,s$ and suppose $(s,q,r)$ is a three-cap.
If $(u,s,q)$ is a cup, append $q$ to obtain a seven-cup. Otherwise
$u$ ends the five-point prefix of the six-cup and starts the four-cap
$(u,s,q,r)$, a forbidden junction. Both possibilities are impossible.
\end{proof}

\begin{lemma}[Anchoring]\label{lem:anchor}
In a heptagon-free set, suppose $(a,b,t_3,t_4,t_5,t_6)$ is a six-cup
and $a$ is a right root. Every six-cup ending at $b$ starts at $a$.
\end{lemma}
\begin{proof}
If a start $x<a$ existed, $(x,a,b)$ would have to be a cap, since
otherwise $x$ could be prepended to the displayed cup. That cap and
the six-cup from $x$ to $b$ would form a heptagon. If $x>a$, the
right root $a$ could be prepended to the six-cup from $x$ to $b$.
Only $x=a$ remains.
\end{proof}

\begin{theorem}[Five-point deletion]\label{thm:deletion}
Every set counted by $D_7$ has a subset of at most five points meeting
every six-cup. In particular $D_7\le B+5$.
\end{theorem}
\begin{proof}
Let $E$ be all actual six-cup endpoints. By \cref{lem:root}, every
triple in $E$ is a cup; thus $|E|\le6$. If $|E|\le5$, delete $E$.
Otherwise write $E=(e_1,\ldots,e_6)$. This is a six-cup, and
\cref{lem:anchor} says that every six-cup ending at $e_2$ starts at
$e_1$. Delete $E\setminus\set{e_2}$. A surviving six-cup would have
an original endpoint in $E$, hence endpoint $e_2$, but would have to
start at the deleted point $e_1$. This is impossible.
The remaining set has no five-cap, six-cup, or heptagon, and has size
at most $B$.
\end{proof}

\section{Finite states and the balanced heptagon exclusion}\label{sec:balanced}
Down-set encodings of monotone paths are classical~\cite{MS,BB}.

The certificate conventions in this section use \emph{cup first, cap
second}. This differs from the order of the arguments of $M_n(a,b)$.
For reference, the two stored conventions are
\begin{center}
\begin{tabular}{ll}
\toprule
Certificate & Coordinates of a cell\\
\midrule
Section~\ref{sec:balanced}: heptagon & (cup, cap)\\
Sections~\ref{sec:octagons}--\ref{sec:319}: octagon & (cap, cup)\\
\bottomrule
\end{tabular}
\end{center}
For $u<v$, let $d(u,v)=(r(u,v),b(u,v))$, where $r+1$ and $b+1$ are
the longest cup and cap lengths ending at the exact pair $(u,v)$.
Define
\begin{equation}\label{eq:states}
 F(v)=\down\set{d(u,v):u<v}.
\end{equation}
All these ranks are actual maxima in the stated ambient set.

\begin{lemma}[Difference and chronology]\label{lem:difference}
For $u<v$, $d(u,v)\in F(v)\setminus F(u)$. Consequently states are
injective, and strict inclusion of two occupied states forces the same
order on their owners.
\end{lemma}
\begin{proof}
If $d(u,v)\in F(u)$, some actual predecessor $t<u$ has
$d(t,u)\ge d(u,v)$ coordinatewise. According to the orientation of
$(t,u,v)$, extending an attaining cup or cap contradicts one of the
two maxima at $(u,v)$. Thus the difference is nonempty. Equal states,
or an earlier state containing a later one, would contradict this
conclusion.
\end{proof}

With no $a$-cap or $b$-cup, the states lie in $[b-2]\times[a-2]$,
whose number of down-sets is $\binom{a+b-4}{a-2}$, including the empty
state. A Ferrers tuple denotes column heights, with trailing zeros
where required by the data format. The finite models permit every
chronological order compatible with inclusion; they do not sort
incomparable states into a preferred order.

\begin{theorem}[A sharp ETV case; computer-assisted]\label{lem:B25}
Every set of 26 points in $\mathbb R^2$, in general position and with
distinct $x$-coordinates, contains a five-cap, a six-cup, or seven
points in convex position. There is a 25-point set avoiding all three.
Thus $M_7(5,6)=25$; reflection also gives $M_7(6,5)=25$.
\end{theorem}
\begin{proof}[Certificate reduction]
For the lower bound, the Erd\H{o}s--Szekeres construction~\cite{ES61},
in the restricted form of Erd\H{o}s--Tuza--Valtr~\cite{ETV}, gives
$\sum_{i=1}^{3}\binom5i=25$ points with none of the three configurations;
see also Baek--Balko~\cite[Section~4]{BB}.
The archive also contains an exact rational 25-point witness and a
separate integer checker confirming maximum cup length five, maximum
cap length four, and no convex heptagon among all
$\binom{25}{7}=480\,700$ seven-point subsets.
This lower-bound check is separate from the Lean upper-bound theorem.

A hypothetical 26-point counterexample has 26 distinct actual states
in $[4]\times[3]$, a universe of 35 states, and hence nine missing
states. The empty state and the one-cell state occur at the first two
points; the third has one of the two two-cell states. In the certificate's
area-sorted dictionary, an eligible hole set therefore avoids indices
0 and 1 and does not contain both 2 and 3. There are exactly
\[
 \binom{33}{9}-\binom{31}{7}=35\,937\,525
\]
eligible hole sets.

The certificate derives necessary hole clauses from pair-label
differences, exact corner attainment, shared predecessor witnesses, actual
owner order, four-point sign constraints, and common-endpoint paths.
A five-cup and four-cap with common endpoints give a heptagon, so
excluding just this type is a necessary restriction. The Lean coverage stage refutes a CNF consisting of 91,720 derived
hole clauses, a cardinality encoding, and one padding tautology:
461 variables and 93,169 clauses in total. Its LRAT proof has
381,647 additions. Separately, two enumerations, recursive and
fixed-popcount, verify that none of the eligible hole sets survives.
\Cref{app:restricted} describes the primitive semantics and proof format.
\end{proof}

Consider now a balanced avoiding set $S$ of size 50. Let $\cupS$ and
$\capS$ be its actual five-cup and five-cap starters, and $\cupE$ and
$\capE$ the corresponding endpoint sets.

\begin{lemma}[Equality partition]\label{lem:equality}
There is an integer $t$ such that
\[
 S=\cupS\sqcup\capE=\capS\sqcup\cupE,
 \qquad |\cupS|=|\capS|=|\cupE|=|\capE|=25,
\]
and, with
\[
 L=\cupS\cap\capS,\quad X=\cupS\cap\cupE,
 \quad Y=\capS\cap\capE,\quad R=\cupE\cap\capE,
\]
one has $|X|=|Y|=t$ and $|L|=|R|=25-t$.
\end{lemma}
\begin{proof}
The set $\cupS$ contains no five-cap: joining the end of such a cap
to an actual outgoing five-cup extends one of them to a six-chain.
Its complement contains no five-cup. By \cref{lem:B25} and reflection,
both parts have size at most 25, and hence exactly 25. Color exchange
and order reversal give the same endpoint counts. A point ending a
five-cap cannot start a five-cup, by the same junction argument, so
the disjoint pairs each exhaust $S$. Taking their four intersections
gives the stated cardinalities.
\end{proof}

\begin{lemma}[Population certificates]\label{lem:populations}
In a planar set with no six-cup, six-cap, or heptagon:
\begin{enumerate}[label=(\roman*)]
\item at most thirteen points both end and start five-cups;
\item any subset avoiding five-cups and five-caps, whose last point
starts a five-cup in the ambient set, has at most nineteen points;
\item if $N_0$ consists of points ending neither five-chain and $N_2$
of points ending both, then
\[\min(|N_0|,|N_2|)\le16.\]
\end{enumerate}
\end{lemma}
\begin{proof}[Necessary reductions and finite certificates]
For (i), select fourteen supposed through-cup points, a five-cup ending
at the first, and a five-cup starting at the last. Exactly 22 points
are selected: fourteen core points and four on each side. Any core
four-cap extends at both ends to a six-cap, so the core has cap length
at most three. Its cap paths extend by one at each end, whereas its
cup paths remain valid ambient paths. The source records exact core
recurrences and these valid implications for ambient ranks, without
requiring ambient attainment in the selection. Its refutation has
3,009 additions.

For (ii), select a supposed 20-point core and the four other vertices
of an actual five-cup from its last point. Incoming ranks ending in
the core are recomputed there; neither color reaches five. The resulting
24-point source is refuted in 79,173 additions.

For (iii), the incoming states of $N_0$ are the twenty ideals of
$[3]^2$. Those of $N_2$ are the twenty ideals
$(4,1+h_1,1+h_2,1+h_3)$ with $h$ an ideal of $[3]^2$. Select seventeen
owners from each family, retaining the ambient first two points.
A named-owner source permits all such selections and all their
inclusion-compatible orders. Pair ranks are ambient and paths are
selected. Its refutation has 102,563 additions.

All three formulas include the necessary four-point packet constraints
and common-endpoint polygon exclusions. The compact formulas and refutations are embedded in the Lean sources.
Their geometric interpretations and logical soundness are proved there;
\cref{sec:formal,app:cnf} describes the interface. Thus each
supposed configuration supplies an assignment to a refuted formula.
\end{proof}

\begin{corollary}\label{cor:range}
In \cref{lem:equality}, $9\le t\le13$. Every twenty ambient five-cup
starters contain an internal five-cup, and a 25-point starter class
has at least six actual internal five-cup sources.
\end{corollary}
\begin{proof}
The through-cup bound gives $t\le13$. In the equality partition,
$N_0=L$ and $N_2=R$, so $25-t\le16$ and $t\ge9$.
Any set of twenty cup starters has no internal five-cap by the
junction argument. If it had no five-cup either, part (ii) of
\cref{lem:populations}, applied to its last point, would be contradicted.
Finally, deleting at most five internal sources from a 25-point starter
class would leave at least twenty ambient starters with no internal
five-cup. Deletion cannot create a new path, so at least six sources
must exist.
\end{proof}

\begin{proposition}[Balanced heptagon bound]\label{prop:balanced}
$C_7\le49$.
\end{proposition}
\begin{proof}
Assume $|S|=50$ and let $K=\cupS$. Retain the 25 points of $K$ and
the four other points of an actual five-cup from $\max K$. For a point
of $K$, no ambient five-cap can end there, so its actual incoming
state lies in the 35-state universe $[4]\times[3]$. Fifteen of these
states have width four; their active population is precisely $t$.
The low-width owners are also cap starters by \cref{lem:equality}.
All 25 states are actual ambient states, not states recomputed after
selecting these 29 points.

Reflect the whole configuration if necessary to make its first triple
a cup, and select the corresponding starter class anew. The four
possible initial packets, in order $012,013,023,123$, are
$1000,1100,1110,1111$. A cup and cap share at most two vertices.
The initial three-cup therefore forces the first five ambient points
to have no ending five-cap, and hence to belong to $K$. In the $1111$
sector the first six points similarly belong to $K$. Exhaustive
literal enumeration retains all seven five-point extensions of $1000$
and all 118 six-point extensions of $1111$ without a six-cup.
The $1100$ and $1110$ sectors remain in the primary formula until
their exclusions are derived.

The primary formula has 151,801 variables and 4,017,399 clauses. It includes \cref{cor:range}, the six-source condition,
free order for incomparable owners, shared ambient pair labels,
selected paths, exact core outgoing paths, and the prefix alternatives.
Exact incoming attainment is imposed only at retained ambient prefix
points, where every predecessor is present. The variable meanings
and necessary clause families are listed in Appendix~\ref{app:cnf}.

A single proof over this unaugmented primary formula derives the empty
clause. It contains 1,989,465 additions and 84,721,022 ordered hints.
Its derivation includes the exclusions of $t=13,9,12$, the two mixed
prefixes, and both final cases $t=10,11$. These case facts are derived
within the proof; the final contradiction has no case assumption.
The compact source uses 190,932 indexed clauses of the primary formula;
the same number of proof additions suffices. Checked replay of
this compact proof establishes unsatisfiability. The formal proof additionally checks each retained clause against a
proved necessary schema under one common assignment.

Every supposed planar $S$ supplies a truthful assignment to the
retained source, so unsatisfiability is a contradiction. Any larger
avoiding set would contain an avoiding fifty-point subset.
\end{proof}

Combining \cref{lem:B25,thm:deletion,prop:balanced,thm:transfer} gives
\[
 \ES(7)\le49+(25+5)+2=81.
\]
The classical construction~\cite{ES61} supplies the lower endpoint
$\ES(7)\ge33$; it is separate from the formalized upper bound.
